%! AP Statistics — Unit 6, Lesson 6.1 — Lesson Plan
\documentclass[10pt]{article}
\usepackage{apstats-article}
\usepackage{apstats-boxes}

\newcommand{\UnitNumberName}{Unit 6: Inference for Categorical Data: Proportions \quad}
\newcommand{\LessonNumberName}{Lesson 6.1: Introducing Statistics --- Why Be Normal?}

\begin{document}
\begin{center}
    {\Large\bfseries \CourseName: \SchoolYear \\
    {\small\bfseries \UnitNumberName \LessonNumberName}}
\end{center}

\begin{tcolorbox}[colback=sky, colframe=navy, boxrule=0.9pt, arc=2mm,
    left=2mm, right=2mm, top=1.5mm, bottom=1.5mm]
    \textbf{Primary Objective:} Students will recognize that distributions of sample statistics
    can vary in shape across repeated samples from the same population, and frame the central
    question of inference: when should we trust that a pattern in data reflects a true population
    characteristic rather than random chance? (Big Idea: VAR)
\end{tcolorbox}

\renewcommand{\arraystretch}{1.6}
\begin{skillbox}[Priority Ideas \& Skills]{goldbox}
\begin{minipage}[t]{0.48\linewidth}
    \textbf{AP Skill 1 --- Selecting Statistical Methods}
    \begin{itemize}
        \item Identify questions suggested by variation in shapes of distributions
              of samples taken from the same population (Skill 1.A; VAR-1.H).
        \item Bridge: Unit 5 built sampling distributions --- Unit 6 uses them to
              make formal decisions about population parameters.
        \item Recognize that shape variation may be random (due to chance) or
              non-random (due to a real population effect): VAR-1.H.1.
    \end{itemize}
\end{minipage}
\hfill
\begin{minipage}[t]{0.48\linewidth}
    \textbf{Key Understandings}
    \begin{itemize}
        \item The question of inference: ``Is the pattern I see in my sample strong
              enough to be convincing evidence about the population?''
        \item Confidence intervals answer: ``Where is the parameter likely to be?''
        \item Significance tests answer: ``Is a claimed parameter value consistent
              with my data?''
        \item This lesson is a framing lesson (not assessed in the CED) that sets
              the stage for Lessons 6.2--6.11.
        \item Normal distributions will be used to model sampling distributions ---
              hence the title ``Why Be Normal?''
    \end{itemize}
\end{minipage}
\end{skillbox}

\begin{skillbox}[{Vocabulary, Concepts, \& Theorems}]{greenbox}
\begin{minipage}[t]{0.48\linewidth}
    \begin{itemize}
        \item \textbf{Statistical inference} --- using sample data to draw conclusions
              about a population parameter
        \item \textbf{Confidence interval} --- a range of plausible values for an
              unknown parameter
        \item \textbf{Significance test} --- a procedure for deciding whether data
              provide sufficient evidence against a claim about a parameter
    \end{itemize}
\end{minipage}
\hfill
\begin{minipage}[t]{0.48\linewidth}
    \begin{itemize}
        \item \textbf{Null hypothesis ($H_0$)} --- the default claim about a parameter
        \item \textbf{Alternative hypothesis ($H_a$)} --- the claim we seek evidence for
        \item \textbf{$p$-value} --- the probability of observing results at least as
              extreme as those observed, assuming $H_0$ is true
        \item \textbf{Normal model} --- the key distribution for inference on proportions
              and means (from Unit 5)
    \end{itemize}
\end{minipage}
\end{skillbox}

\begin{skillbox}[Activate Prior Knowledge \& Spiral Review]{sky}
\begin{minipage}[t]{0.48\linewidth}
    \textbf{Prior Knowledge Check (3--4 min)}
    \begin{itemize}
        \item Unit 5: Sampling distribution of $\hat p$ --- shape, center $\mu_{\hat p} = p$,
              spread $\sigma_{\hat p} = \sqrt{p(1-p)/n}$.
        \item Unit 5: Large Counts condition ($np \ge 10$, $n(1-p) \ge 10$) ensures
              the sampling distribution is approximately normal.
        \item Connect: We used sampling distributions to compute probabilities of outcomes
              from known populations. Now the population is \emph{unknown} --- we flip the question.
    \end{itemize}
\end{minipage}
\hfill
\begin{minipage}[t]{0.48\linewidth}
    \textbf{Connections Forward}
    \begin{itemize}
        \item Confidence intervals for one proportion $\to$ Lessons 6.2--6.3
        \item Significance tests for one proportion $\to$ Lessons 6.4--6.6
        \item Type I/II errors and power $\to$ Lesson 6.7
        \item Two-sample procedures $\to$ Lessons 6.8--6.11
        \item Inference for means $\to$ Unit 7
    \end{itemize}
\end{minipage}
\end{skillbox}

\begin{skillbox}[Hook \& Lesson]{sky}
\begin{minipage}[t]{0.48\linewidth}
    \textbf{Hook (5--7 min)}\\
    Display: A news headline reads ``Survey: 53\% of voters support Candidate A.''
    A second headline the same day reads: ``New poll: only 47\% back Candidate A.''\\[4pt]
    Ask: ``Can both be right? What should we believe?''\\[4pt]
    Reveal: Both are plausible outcomes of random sampling from the same electorate.
    The questions of Unit 6: \textbf{How do we use sample data to estimate the true
    proportion? How do we test whether a claim about the true proportion is supported?}
\end{minipage}
\hfill
\begin{minipage}[t]{0.48\linewidth}
    \textbf{Lesson Flow (15 min)}
    \begin{enumerate}
        \item Revisit the sampling distribution of $\hat p$ from Unit 5: shape is
              approximately normal when Large Counts holds; center is $p$; SD is
              $\sqrt{p(1-p)/n}$.
        \item The key idea: if $\hat p$ varies predictably around $p$, we can
              \emph{work backwards} from $\hat p$ to estimate $p$, or to test a
              claimed value of $p$.
        \item Preview the two big inference tools: a \textbf{confidence interval}
              captures $p$ with a stated level of certainty; a \textbf{significance
              test} evaluates evidence against a specific value of $p$.
        \item Briefly show how both tools rely on the normal model for $\hat p$.
    \end{enumerate}
\end{minipage}
\end{skillbox}

\begin{skillbox}[Explicit Instruction \& Active Monitoring]{sky}
\begin{minipage}[t]{0.48\linewidth}
    \textbf{Direct Instruction (8 min)}
    \begin{itemize}
        \item Compare the Unit 5 question (``Given $p = 0.6$, what is the probability
              that $\hat p > 0.65$?'') with the Unit 6 question (``Given $\hat p = 0.65$,
              what can we say about $p$?'').
        \item Introduce the ``inversion'' idea: inference flips probability from forward
              (known $p$, unknown $\hat p$) to backward (known $\hat p$, unknown $p$).
        \item Emphasize: the normal model still does the work because the sampling
              distribution is approximately normal when conditions are met.
    \end{itemize}
\end{minipage}
\hfill
\begin{minipage}[t]{0.48\linewidth}
    \textbf{Active Monitoring}
    \begin{itemize}
        \item Listen for students confusing a confidence interval with a sampling
              distribution probability. Redirect: ``A CI is a range for the parameter,
              not for the statistic.''
        \item Cold-call: ``What conditions must hold before we can use a normal model
              for $\hat p$?''
        \item Check warm-up for fluency with the Large Counts condition and $SE_{\hat p}$.
    \end{itemize}
\end{minipage}
\end{skillbox}

\begin{skillbox}[Group Work \& Differentiation]{redbox}
\begin{minipage}[t]{0.48\linewidth}
    \textbf{Partner Activity (8--10 min)}\\
    Three survey scenarios with different sample proportions and sample sizes.
    For each:
    \begin{enumerate}
        \item Verify Large Counts holds.
        \item Compute $SE_{\hat p}$.
        \item Sketch the approximate sampling distribution.
        \item Identify what inference tool you would use (CI or test) to answer
              the research question.
    \end{enumerate}
\end{minipage}
\hfill
\begin{minipage}[t]{0.48\linewidth}
    \textbf{Differentiation}
    \begin{itemize}
        \item \textit{Support (Tier R)}: Provide the formula for $SE_{\hat p}$ and a
              completed example. Ask students to identify the inference goal only.
        \item \textit{On-grade (Tier A)}: Complete all four steps for each scenario.
        \item \textit{Extension (Tier E)}: For one scenario, describe how the interval
              would change if $n$ were quadrupled. How does that change $SE_{\hat p}$?
        \item \textit{ELL}: Bilingual card --- ``confidence interval / intervalo de confianza,''
              ``significance test / prueba de significancia.''
    \end{itemize}
\end{minipage}
\end{skillbox}

\begin{skillbox}[Individual Work \& Assessment]{redbox}
\begin{minipage}[t]{0.48\linewidth}
    \textbf{Exit Ticket (5 min)}
    \begin{enumerate}
        \item A random sample of 200 students finds $\hat p = 0.64$ use school
              transportation. State the research question as an inference question
              (estimate or test?).
        \item Verify the Large Counts condition for this sample.
        \item Calculate $SE_{\hat p}$. Show your work.
    \end{enumerate}
\end{minipage}
\hfill
\begin{minipage}[t]{0.48\linewidth}
    \textbf{AP-Style Check}\\
    \textit{A researcher collects a random sample and computes $\hat p = 0.42$. She wants
    to decide whether the population proportion is higher than 0.35. The appropriate
    next step is to:}
    \begin{enumerate}[label=(\Alph*)]
        \item Compute the population proportion.
        \item Perform a significance test.
        \item Compute a confidence interval for $\hat p$.
        \item Re-sample with a larger $n$.
    \end{enumerate}
    Answer: (B)
\end{minipage}
\end{skillbox}

\begin{skillbox}[Reinforcement \& Extension]{goldbox}
\begin{minipage}[t]{0.48\linewidth}
    \textbf{Homework / Practice}
    \begin{itemize}
        \item For three given scenarios (each with a context, $n$, and $\hat p$):
              (a) verify Large Counts; (b) compute $SE_{\hat p}$; (c) classify the
              research question as estimation or hypothesis testing.
        \item Reflection: ``Why is the normal model essential for the inference
              procedures we'll study in Unit 6?''
    \end{itemize}
\end{minipage}
\hfill
\begin{minipage}[t]{0.48\linewidth}
    \textbf{Extension / Preview}
    \begin{itemize}
        \item \textit{Extension}: Research a published poll. Identify the sample proportion,
              sample size, and state whether the researchers appear to have reported a CI
              or conducted a significance test.
        \item \textit{Preview}: Lesson 6.2 introduces the one-sample $z$-interval for
              a population proportion --- the first complete inference procedure.
    \end{itemize}
\end{minipage}
\end{skillbox}

\end{document}
